\chapter{The Mirror: The West Runs the Playbook Against Itself}
\label{c:mirror}

\begin{glance}
Over 2026 the largest Western firms in the sector converged on the same method---not in imitation, and in most cases without reference to China at all, but by the same economic logic arriving at the same instrument. Each now commoditizes, deliberately and at its own near-term expense, a layer it does not own, to defend the layer it does: the accelerator vendor gives away models and the training data with them, the laboratories design accelerators, the hyperscalers rebuild the scale-up fabric out of merchant parts, and a launch company is building a fab. The strategy this book has attributed to Beijing is no longer Chinese; what remains asymmetric is not the method but the scope. Two consequences follow in opposite directions, so both are printed: Proposition~P3 acquires a second mechanism that would not require a Chinese actor to arrive [S, contingent], and a West that commoditizes its own layers competently may capture the surplus of a cheap-intelligence world as surely as a China that does. \full{17}
\end{glance}

\section{Who commoditizes what}

Nothing in \S\ref{c:theory}'s account of openness is national. A firm holding layer~$A$ and facing a rent-taker at the adjacent layer~$B$ can drive $B$ to marginal cost and collect the transferred surplus at $A$---losing nothing it was earning, expanding its market as the complement cheapens, and imposing the loss on a firm it does not own. What is new in 2026 is the simultaneity and the candour: four distinct classes of Western actor executed the move at scale within a single year, three of them on the record (Table~\ref{tab:mirror}). The question it provokes is the one this chapter exists to answer---when everyone is giving away everything, where does the rent actually land?

\begin{table}[htbp]
\centering\footnotesize
\caption{Who commoditizes what, as of \datafreeze. ``Instrument'' names the vehicle, not the intent.}
\label{tab:mirror}
\begin{tabular}{@{}p{20mm}p{28mm}p{28mm}p{46mm}p{14mm}@{}}
\toprule
\textbf{Actor} & \textbf{Layer defended} & \textbf{Layer commoditized} & \textbf{Instrument} & \textbf{Status} \\
\midrule
China & Hardware volume, domestic cloud, standards & \emph{All of them} & Open weights, harness, ISA, interconnect; mandate; power tariff & Running \\
NVIDIA & Accelerator and CUDA & The model; the financing of the datacentre & Nemotron weights \emph{and} training data; cuTile (Apache~2.0); Poolside ``Model Factory'' licence and 109 researchers; residual value guarantees & Partial (revenue share paused) \\
AMD & Accelerator & The model; the scale-up fabric & Instella (weights, data, code); ROCm; UALink over Ethernet; customer warrants & Running \\
OpenAI & Frontier model and platform & The accelerator and the system & Jalape\~no with Broadcom; 10\,GW term sheet & Running \\
Google, Microsoft, Meta, AWS & Cloud platform and applications & The accelerator, system and fabric & TPUs sold into customer sites; Maia~200; MTIA; Trainium; Ethernet scale-up & Running \\
Broadcom, Marvell & Co-design and networking & The merchant GPU & Custom ASIC programmes; eRoCE & Running \\
SpaceX, Tesla & Launch, vehicles, robots & The fab & Terafab; dual-foundry AI5/AI6 & Announced \\
\bottomrule
\end{tabular}
\end{table}

\section{The accelerator vendor gives away the model---and underwrites the buildout}

The clearest instance is the largest company in the sector, and it runs the move on two layers at once. \emph{The model layer.} NVIDIA released the Nemotron~3 family with not only open weights but the \emph{training datasets and the reinforcement-learning libraries}, and Nemotron~3.5 Lightning followed in August 2026 under the OpenMDW-1.1 licence with data and recipes, atop an open pretraining collection exceeding ten trillion tokens~\cite{nemotron}; TensorRT-LLM is Apache~2.0, and the CUDA Tile IR and cuTile compiler stack were open-sourced under the same licence~\cite{cuda-tile}. The position is measurable rather than rhetorical: in the week to 29~August 2026 Nemotron~3 Ultra was the sixth-largest model by tokens routed on the largest neutral router---on its \emph{free} endpoint, and router volumes are a small and promotion-distorted slice of global inference~\cite{openrouter-aug26} [A]---a chip company is a top-ten model provider by volume. The rationale is on the record: ``nearly all open models run on NVIDIA\ldots\ our position in open models is very good, because the CUDA ecosystem is literally everywhere,'' and, from the same firm's director of developer technology, that durable value accrues to proprietary \emph{data} rather than to weights---which is exactly why it can release the weights and the data together~\cite{nvda-open-models} [P]. That is the argument this book attributed to Chinese laboratories, made in public by the incumbent whose rent the Chinese version of the move is aimed at.

\emph{And in August 2026 the firm stopped merely publishing weights and began buying the means of producing them.} NVIDIA and Poolside disclosed a transaction of approximately \$7 billion: \$6 billion for a \emph{non-exclusive} licence to Poolside's ``Model Factory'' training platform, \$1 billion of equity, and offers to \textbf{109 Poolside engineers and researchers} to join NVIDIA and work on Nemotron, the stated purpose being a leading American open-weight model able to compete with DeepSeek and Moonshot~\cite{nvda-poolside} [P]. On 26~August \emph{The Information} reported that NVIDIA was closing on an acquisition of Hugging Face for approximately \$12.9 billion; at \datafreeze\ neither party had confirmed it and other reporting held that no agreement was signed, so this book does not assert that it happened and what follows is written to survive its not happening~\cite{nvda-hf} [A]. The analytical point does not need the second item. A firm is acquiring the \emph{supply chain} of the open layer---the training factory, the staff who operate it, and, if the report confirms, the distribution hub through which very nearly all open weights reach very nearly all users---and a supply chain under one owner is a different object from a commons, however permissive the licences on what comes out of it. The form is a pattern: roughly \$27 billion across Groq, Enfabrica and Poolside in nine months, each structured so the target stays nominally independent and outside merger review, with a letter from Senators Warren and Blumenthal saying NVIDIA has ``effectively acquired Groq in all but name'' and no enforcement action following by the currency date~\cite{nvda-structured} [V/P]. An equity acquisition of the hub \emph{would} require review, which is why the unconfirmed report matters more than its price.

\emph{The datacentre layer.} The second move is larger and gives the first its durability: through 2026 NVIDIA converted AI compute from a capital good its customers had to finance into an asset class third-party capital would finance, by standing behind the residual value. The filed positions supersede much looser reporting: residual value guaranties on datacentre leases at Pike County, Ohio, whose tenant is an OpenAI affiliate---roughly 4.25\,GW of IT load, a twenty-year maximum term, an aggregate obligation ``cumulatively capped at \$105 billion''~\cite{nvda-ohio} [P], with Chapter~\ref{c:trade} pricing the wider exposure class---and memoranda with six asset managers to establish platforms mobilizing ``over \$500 billion of third-party capital''~\cite{nvda-500b} [P]. The regulatory precondition is almost never reported alongside it: a fortnight earlier the SEC's Office of Structured Finance took the position that securities in datacentre securitizations are not asset-backed securities under the Exchange Act, placing them outside the five per cent risk-retention rule~\cite{sec-dcs} [V]. Read them together: if the model layer is a commons anyone can serve, and the datacentre layer is financed against the residual value of the accelerator rather than the creditworthiness of any tenant, the failure of any individual model company does not remove demand for the silicon. The commons survives its sponsors.

\emph{The counter-evidence is substantial and belongs first in any honest summary.} The programme was partly suspended: on 27~August NVIDIA paused parts of the July revenue-sharing and credit-support scheme, reportedly after employees flagged potential antitrust exposure~\cite{nvda-pause} [A], and a programme halted within eight weeks of launch is not the execution of a long-laid plan. The trajectory is retrenchment rather than escalation---a \$100B commitment to OpenAI became a \$30B equity position, and talks reported at up to \$250B for the Ohio structure became \$105B on the filing---and the structures are narrower than the headline, since the Ohio guarantee \emph{terminates early if OpenAI achieves a satisfactory credit rating} and OpenAI must reimburse and indemnify NVIDIA for every amount paid~\cite{nvda-ohio}: a bridge to creditworthiness rather than permanent underwriting, and one management denies is vendor financing at all~\cite{nvda-critique}. Against all of that stands one line in the same filings, which supports the circular reading better than any critic's argument and which Chapter~\ref{c:trade} takes apart: GAAP earnings per share exceeded the non-GAAP figure, because other income was almost entirely gains on an equity portfolio concentrated in the firms that buy the chips~\cite{nvda-q2fy27} [P].

AMD runs the same playbook, and that matters for the diagnosis: Instella released \emph{fully} open---weights, data and training code---ROCm against CUDA, the UALink consortium building the scale-up fabric over merchant Ethernet~\cite{amd-open}, and warrants for up to 160 million shares at one cent issued to OpenAI and again to Meta, vesting against gigawatt purchase tranches, alongside \$4.1B of disclosed datacentre lease guarantees~\cite{amd-warrants,amd-q2-2026} [P]. Stated plainly, the customer is being paid in the vendor's own equity to buy the vendor's product. That two rivals independently adopt so unusual an instrument argues for competitive necessity at least as strongly as for grand design, and the book holds both readings.

\section{The other direction: the laboratories and the hyperscalers commoditize the silicon}

Hot Chips 2026 is where this stopped being describable as a hedge~\cite{hc2026}. OpenAI presented Jalape\~no, an inference accelerator co-designed with Broadcom---13.4 PFLOP/s at four bits, 216\,GiB of HBM4---under a term sheet covering ten gigawatts of deployment through 2029~\cite{hc2026-jalapeno,openai-broadcom} [P]. Three weeks later there were numbers, and they are the sharpest claim any custom part has made against the merchant GPU: 1.5--1.9$\times$ more AI work per watt at peak throughput and 1.7--3.6$\times$ lower end-to-end latency against GB200, GB300 and MI355X systems~\cite{jalapeno-results} [P]. Two qualifications are load-bearing rather than decorative---\emph{the part was not deployed at the currency date}, so every figure is a vendor-run measurement of a part in no production service, and the comparison sets HBM4 parts against older HBM3E systems, which SemiAnalysis calls ``somewhat incomplete and unfair.'' And the same post carries the sentence that constrains every reading of it: ``We will continue to widely deploy accelerators from NVIDIA and other partners for both training and inference workloads''~\cite{jalapeno-results} [P]. Set beside Huang's statement that OpenAI's existing and planned commitments represent approximately twelve gigawatts of NVIDIA compute through 2030~\cite{nvda-huang-q2}, the two firms are saying the same thing from opposite directions: a layer being commoditized by a firm that goes on buying from it. That is not the shape of displacement. Google will begin delivering TPUs to selected customers \emph{in those customers' own datacentres}~\cite{hc2026-tpu8,tpu-external} [P]; Microsoft's Maia~200 makes an economic rather than an architectural claim, thirty per cent better performance per dollar than the latest generation in its fleet~\cite{maia200}; Meta has four MTIA generations on a six-month cadence~\cite{mtia400}; AWS states more than a million Trainium processors deployed~\cite{trainium3}; and behind all of them Broadcom's AI semiconductor revenue grew 143 per cent to \$10.8B in a quarter with the next guided above \$16B~\cite{broadcom-q2fy26} [P]. Annualized against NVIDIA's datacentre run rate of \$89.0B in the quarter that places merchant custom-ASIC revenue at roughly eleven to fifteen per cent of the combined total [M]: not yet a substitute for the merchant GPU, but growing several times faster and bought by precisely the firms whose purchase orders set NVIDIA's revenue. The fabric tells the same story more sharply, because there the substitution has already happened: of eight parts at Hot Chips with an off-package scale-up fabric, six build it from Ethernet-class merchant silicon and NVIDIA alone retains a proprietary one~\cite{hc2026-scaleup}. The layer that was the deepest source of lock-in three years ago is being rebuilt out of parts anyone can buy.

The counter-evidence here is decisive enough to print immediately. On 26~August 2026 AWS and NVIDIA announced two million additional NVIDIA GPUs across 2027--28 and, materially, that \emph{Trainium4 will adopt NVIDIA's NVLink Fusion interconnect and custom HBM}~\cite{aws-nvda} [P]: the largest deployer of non-NVIDIA AI silicon in the world is placing its own accelerator inside NVIDIA's rack-scale interconnect and memory stack. Set against a quarter in which datacentre revenue rose 117 per cent to \$89.0B, against total-company third-quarter guidance of \$108.0B~\cite{nvda-q2fy27}, the honest reading is not that the incumbent is being displaced by the custom-silicon wave but that it is absorbing it: commoditization has so far moved the boundary of what NVIDIA sells rather than the size of it. One correction of the record belongs here, dryly, because a false version circulates: it was not NVIDIA's datacentre head who left for OpenAI. It was \emph{OpenAI's own} Head of Data Centers who left OpenAI in August 2026 after sixteen months---and that title was OpenAI's alone: he arrived as a distinguished engineer of nearly five years at Meta and almost twelve at Google on datacentre infrastructure, and had never worked at NVIDIA; OpenAI's infrastructure organization is led by an executive recruited from \textbf{Intel}; NVIDIA's datacentre leadership is intact; and the senior technical flow of 2026 ran \emph{toward} NVIDIA, which paid roughly \$27 billion for two such flows---a market in which talent must be bought at that price is not obviously a market at rest~\cite{openai-exec,nvda-segments} [V].

At the bottom of the stack, SpaceX confirmed a \$16.8B first phase for ``Terafab,'' a vertically integrated logic and memory site in Texas~\cite{terafab}. This takes the coldest grade in the chapter [R]---the announced figures have moved three times, and no process node or timeline was disclosed---and it is evidence of intent rather than capacity. What it demonstrates is the reach of the logic: an actor whose rent is in launch, vehicles and robots, facing a fabrication layer it regards as a constraint, reasons its way to owning the constraint.

\section{The asymmetry that survives}

Two asymmetries survive the convergence of method, and they run in opposite directions.

\textbf{The first favours the West, and the book should say so plainly.} Commoditization is not losing: a stack in which every layer is cheap is one in which the surplus accrues to whoever holds the layer that stays scarce, and at the currency date those are mostly still Western-held---the CUDA installed base, the enterprise deployment surface, the trust apparatus of Chapter~\ref{c:trust}, and the deployment-generated data that is Chapter~\ref{c:elements}'s fifth element. The clearest demonstration that the West is executing rather than retreating is the price inversion of Chapter~\ref{c:elements}: the incumbent priced the bottom tier of its ladder \emph{below the list price of most Chinese open-weight models} and took roughly three-quarters of the \emph{gain} from competitors~\cite{price-inversion} [P/A]. A firm that can price underneath the commodity and gain share doing it is not a firm being commoditized. It is a firm commoditizing.

\textbf{The second favours China, and it is structural rather than tactical.} Every Western move in Table~\ref{tab:mirror} commoditizes a layer in which \emph{another Western firm} was earning the rent: NVIDIA destroys the model layer's rent to protect silicon, the laboratories and hyperscalers destroy the silicon layer's rent to protect the model and the platform. Each decision is individually rational, and jointly they \emph{redistribute} the Western rent pool rather than necessarily shrinking it---so the pincer claim is conditional, biting to the extent that the layers the rent moves \emph{to} are layers a Chinese actor can also occupy. What China gains is not the destruction of a rent pool but a cheaper entry condition: it needs only to be present in both layers when the rents have moved, and presence is very much cheaper than victory.

Sharper still, and this is the sentence the chapter exists to reach: \emph{the layers the West is commoditizing are disproportionately the layers where China is behind, and several of these moves directly reduce the value of Western state instruments}---open weights shipped with their training data, an open ISA qualified as a CUDA host~\cite{hc2026-nv-riscv}, a scale-up domain rebuilt out of merchant Ethernet~\cite{hc2026-scaleup}. None is a concession to China; all are commercial decisions taken against a domestic rival, and the externality is the same either way. Chapter~\ref{c:strategies} takes the policy consequence, which is that a denial policy premised on the denied party having to obtain the capability from the controlling country is undermined as effectively by an American release as by a Chinese one.

\section{What the mirror licenses, and what it forbids}

It licenses three claims, all of which Chapter~\ref{c:summary} carries forward: that the playbook of Chapter~\ref{c:precedents} is a general theory of a contest over a layered stack rather than a description of Chinese exceptionalism; that the commoditization of the AI stack is \emph{overdetermined}, no longer contingent on any single actor's strategy holding; and that the West's own commercial decisions are eroding its own policy instruments, which is a finding for Chapter~\ref{c:strategies} rather than a flourish. It forbids three equally. That the West is passively losing---it is not. That commoditization implies a Chinese victory---it implies only that the rent moves, and whether it moves to a layer China holds is the open question of the whole book. And that the financing structures are straightforwardly vendor financing---management denies it in terms, the largest structure terminates on a credit-rating event and carries full indemnification, one programme was suspended within eight weeks over antitrust concerns, and the honest verdict at \datafreeze\ is that the practice is unusual, large, growing, disclosed, and not yet demonstrated to be either sound or unsound. One item is held open so a reader can date it: if the reported acquisition of the distribution hub confirms, the description of the open-model layer above changes in kind rather than in degree, because a commons whose warehouse belongs to the vendor of the complement is a distribution channel with a permissive licence attached.
